% block: linear-systems-heading page:14 type:section
\subsection{{线性方程组}}\label{sec:linear-systems}
% block: linear-system-model page:14 type:ordinary-paragraph
{本书考虑下列形式的线性方程组:}
\begin{equation}\label{eq:linear-system}
A\bm x=\bm b,
\end{equation}
{其中 $A\in\mathbb C^{N\times N}$, $\bm b\in\mathbb C$, 并假设系数矩阵 $A$ 非奇异, 即 $A$ 的逆存在.}
% block: norm-motivation page:14 type:ordinary-paragraph
{设 $\tilde{\bm x}$ 为线性方程组 \eqref{eq:linear-system} 的一个近似解. 定量了解解 $\bm x$ 与近似解 $\tilde{\bm x}$ 之间距离的一种方法, 是计算 $\bm x-\tilde{\bm x}$ 的向量范数. 类似地, 两个矩阵 $A,B$ 之间的距离可由 $A-B$ 的矩阵范数定量衡量. 向量范数和矩阵范数的定义将在下面两个小节中介绍. 这里的说明以 \BookCite{170} 为基础; 关于数值线性代数的详细内容, 参见 \BookCite{45}, \BookCite{79}, \BookCite{107}, \BookCite{191}.}
% block: vector-norm-heading page:15 type:subsection
\subsubsection{{向量范数}}\label{subsec:vector-norm}
% block: vector-norm-definition page:15 type:ordinary-paragraph
{向量 $\bm x:=(x_1,x_2)^\top\in\mathbb R^2$ 的长度为 $(x_1^2+x_2^2)^{1/2}$. 范数是长度概念的推广. 如果对所有 $\bm x,\bm y\in\mathbb R^n,\alpha\in\mathbb R$, 映射 $\|\cdot\|:\mathbb R^n\to\mathbb R$ 满足下列规则, 就称它为向量范数:\footnote{一般地, 给定向量空间 $V$, 如果映射 $\|\cdot\|:V\to\mathbb R$ 对所有 $\bm x,\bm y\in V,\alpha\in\mathbb K$ ($\mathbb K=\mathbb R$ 或 $\mathbb C$) 满足 (\ref{cond:vector-norm-positive})--(\ref{cond:vector-norm-triangle}), 就称它为范数.}}
\begin{enumerate}[label=(Nv\arabic*),ref=Nv\arabic*]
\item\label{cond:vector-norm-positive} $\|\bm x\|\ge0;\quad\|\bm x\|=0\iff\bm x=\bm0$.\quad{(正性)}
\item\label{cond:vector-norm-homogeneous} $\|\alpha\bm x\|=|\alpha|\|\bm x\|$.\quad{(齐次性)}
\item\label{cond:vector-norm-triangle} $\|\bm x+\bm y\|\le\|\bm x\|+\|\bm y\|$.\quad{(三角不等式)}
\end{enumerate}
% block: vector-p-norms page:15 type:ordinary-paragraph
{例如, $p$-范数 ($p\ge1$) 定义为}
\[
\|\bm x\|_p:=(|x_1|^p+|x_2|^p+\cdots+|x_n|^p)^{\frac1p}.
\]
{当 $p=1,2,\infty$ 时, 向量 $p$-范数写为}
\begin{itemize}
\item $\|\bm x\|_1=\sum_{i=1}^n|x_i|$,
\item $\|\bm x\|_2=(\sum_{i=1}^n|x_i|^2)^{\frac12}\ (=\sqrt{\bm x^\top\bm x})$,
\item $\|\bm x\|_\infty=\max_{1\le i\le n}|x_i|$.
\end{itemize}
% block: a-norm-convention page:15 type:ordinary-paragraph
{设 $A$ 为对称正定矩阵, 即 $A$ 对称且其全部特征值均为正. 则称 $\|\bm x\|_A:=\bm x^\top A\bm x$ 为 $\bm x$ 的 $A$-范数; Section~\ref{subsec:conjugate-gradient} 将使用这一范数.}
{设 $\tilde{\bm x}$ 为线性方程组 \eqref{eq:linear-system} 的一个近似解. 则 $\|\bm x-\tilde{\bm x}\|_p$ 是 \eqref{eq:linear-system} 的精确解与近似解之间的距离.}
{本书中除非另有说明, 符号 $\|\cdot\|$ 都表示 2-范数, 即 $\|\cdot\|=\|\cdot\|_2$.}
% block: matrix-norm-definition pages:15,16 type:ordinary-paragraph
\subsubsection{{矩阵范数}}\label{subsec:matrix-norm}
{类似于向量范数的定义, 矩阵范数定义如下: 如果映射 $\|\cdot\|:\mathbb R^{m\times n}\to\mathbb R$ 对所有 $A,B\in\mathbb R^{m\times n},\alpha\in\mathbb R$ 满足下列三项, 就称它为矩阵范数:}
\begin{enumerate}[label=(Nm\arabic*),ref=Nm\arabic*]
\item\label{cond:matrix-norm-positive} $\|A\|\ge0;\quad\|A\|=0\iff A=O$,\quad ($O$: {零矩阵})\quad{(正性)}
\item\label{cond:matrix-norm-homogeneous} $\|\alpha A\|=|\alpha|\|A\|$,\quad{(齐次性)}
\item\label{cond:matrix-norm-triangle} $\|A+B\|\le\|A\|+\|B\|$.\quad{(三角不等式)}
\end{enumerate}
{此外, 若矩阵范数满足}
\begin{enumerate}[label=(Nm\arabic*),ref=Nm\arabic*,start=4]
\item\label{cond:matrix-norm-product} $\|AB\|\le\|A\|\|B\|,\quad A\in\mathbb R^{m\times n},B\in\mathbb R^{n\times q}$,
\end{enumerate}
{则称该矩阵范数具有次乘性. 注意, 有些著作也将满足 (\ref{cond:matrix-norm-positive})--(\ref{cond:matrix-norm-product}) 的映射称为矩阵范数. 作为具有次乘性的矩阵范数的一个例子, 下列从属矩阵范数被广泛使用:\footnote{Sup 可以替换为 max.}}
\begin{equation}\label{eq:subordinate-matrix-norm}
\|A\|_p:=\sup_{\bm x\ne\bm0}\frac{\|A\bm x\|_p}{\|\bm x\|_p}.
\end{equation}
% block: matrix-norm-examples page:16 type:ordinary-paragraph
{当 $p=1,2,\infty$ 时, $\|A\|_p$ 如下:}
\begin{itemize}
\item $\|A\|_1=\max_j\sum_{i=1}^n|a_{ij}|$,
\item $\|A\|_2=\sqrt{\text{maximum eigenvalue of }A^\top A}$,
\item $\|A\|_\infty=\max_i\sum_{j=1}^n|a_{ij}|$.
\end{itemize}
{Frobenius 范数也具有次乘性:}
\begin{itemize}\item $\|A\|_{\mathrm F}:=\sqrt{\sum_{i=1}^m\sum_{j=1}^n|a_{ij}|^2}$.\end{itemize}
{设 $\bm a_1,\ldots,\bm a_n$ 为矩阵 $A\in\mathbb R^{m\times n}$ 的列向量, 并令 $\bm a=[\bm a_1^\top,\ldots,\bm a_n^\top]^\top$, 即将全部列首尾相接得到一个长向量 $\bm a$. 则 $A$ 的 Frobenius 范数等于 $\bm a$ 的 2-范数, 即 $\|A\|_{\mathrm F}=\|\bm a\|_2$.}
% block: max-norm page:16 type:ordinary-paragraph
{下列形式的最大值范数}
\[
\|A\|_{\max}:=\max_{i,j}|a_{ij}|
\]
{不具有次乘性. 例如, 考虑所有元素均为 1 的 2 行 2 列矩阵 $A$. 则 $\|AA\|_{\max}=\|A^2\|_{\max}=2$, 而 $\|A\|_{\max}\|A\|_{\max}=1$. 因此, 最大值范数不满足 (\ref{cond:matrix-norm-product}), 因为 $\|AA\|_{\max}>\|A\|_{\max}\|A\|_{\max}$.}
% block: matrix-norm-properties page:16 type:theorem
{下面给出矩阵范数的一些性质.}
\begin{Theorem}\label{thm:matrix-norm-properties}
{下列性质成立:}
\begin{enumerate}[label=(\arabic*),ref=\arabic*]
\item\label{part:matrix-norm-action} $\|A\bm x\|_p\le\|A\|_p\|\bm x\|_p,\quad\|A\bm x\|_2\le\|A\|_{\mathrm F}\|\bm x\|_2$,
\item $\|AB\|_p\le\|A\|_p\|B\|_p,\quad\|AB\|_{\mathrm F}\le\|A\|_{\mathrm F}\|B\|_{\mathrm F}$,
\item $\|QA\tilde Q\|_2=\|A\|_2,\quad\|QA\tilde Q\|_{\mathrm F}=\|A\|_{\mathrm F}$.
\end{enumerate}
{这里, $Q$ 和 $\tilde Q$ 是正交矩阵.\footnote{若实方阵 $Q$ 满足 $Q^\top Q=I$, 则称其为正交矩阵, 其中 $I$ 是单位矩阵.}}
\end{Theorem}
% block: matrix-norm-comparisons page:16 type:theorem
\begin{Theorem}\label{thm:matrix-norm-comparisons}
{设 $A$ 为 $n\times n$ 矩阵. 则}
\begin{enumerate}[label=(\arabic*)]
\item $n^{-1/2}\|A\|_2\le\|A\|_1\le n^{1/2}\|A\|_2$,
\item $n^{-1/2}\|A\|_2\le\|A\|_\infty\le n^{1/2}\|A\|_2$,
\item $n^{-1}\|A\|_\infty\le\|A\|_1\le n\|A\|_\infty$,
\item $\|A\|_1\le\|A\|_{\mathrm F}\le n^{1/2}\|A\|_2$.
\end{enumerate}
\end{Theorem}
